\scanpage{48}
% L1–4
Soit $x\unsure{\mathord{-}|\to}y$ un mono, considérons $x'',y''$ les plus grands sous-objets de $x,y$ $\in\mathrm{Ob}\,C''$. Il faut prouver que
\diagram{\begin{tikzcd}[ampersand replacement=\&,row sep=0.1em,column sep=0.3em]
x \&[-1em] \subset \&[-1em] y\\
\cup \& \& \cup\\
x'' \& \subset \& y''
\end{tikzcd}}
$x''=x\cap y''$, or c'est trivial.

% L5
\uline{Cor} $\underline{\Gamma}'(x)/\underline{\Gamma}''(x)\hookrightarrow R^0(\underline{\Gamma}'/\underline{\Gamma}'')(x)$

% L6–7
Considérons pour $x\in\mathrm{Ob}\,C$ variable
\[
\underline{\Gamma}'x\longrightarrow\Lambda\underline{\Gamma}'x
\]

% L8–9
d'où un hom
\[
\underline{\Gamma}'/\underline{\Gamma}''(x)\longrightarrow\Lambda\underline{\Gamma}'(x)
\]

% L10–12
et comme $\Lambda\underline{\Gamma}'$ est exact à g.
\diagram{\begin{tikzcd}[ampersand replacement=\&]
\underline{\Gamma}'x \arrow[d] \arrow[dr] \& \\
R^0\underline{\Gamma}'/\underline{\Gamma}'' \arrow[r] \& \Lambda\Gamma'
\end{tikzcd}}

% L13–14
\uline{Prop} C'est là un \uline{isom} (?)

équivaut :

% L13–17, left margin
% slanted
\margin{\uline{Faux en} \uline{général} ! \uline{Ok. si les} \uline{objets injectifs} de \uline{$C''$ sont} \uline{injectifs dans $C$}}

% L15–17
\uline{Cor} Pour tt $x\in\mathrm{Ob}\,C$, on a ceci

a) $R^0(\underline{\Gamma}'/\underline{\Gamma}'')(x)$ est $C''$-\uline{fermé} [dans $\underline C'$, ou dans $\underline C$ c'est pareil...]

b) $\underline{\Gamma}'x/\underline{\Gamma}''(x)\longrightarrow R^0(\underline{\Gamma}'/\underline{\Gamma}'')(x)$ est un $\underline C''$-isom.

% L18–19
Cela équivaut aussi (c'est \unsure{animé} par un \unsure{Mons|Mais} \unsure{la|le} \deleted{\unsure{vérif}} \unsure{vérifier})

% L20–21
\uline{Cor} Si $x$ est injectif, alors $\underline{\Gamma}'(x)/\underline{\Gamma}''(x)\xrightarrow{\sim}\Lambda\underline{\Gamma}'(x)$, i.e. $\underline{\Gamma}'(x)/\underline{\Gamma}''(x)$ est \deleted{$\unsure{\underline{\Gamma}'|\underline C'}$} $\underline C''$-fermé.

% L22–23
En effet, $x/\underline{\Gamma}''(x)$ est $C''$-fermé, \deleted{et on voit} \inserted{[en effet, il est injectif par loc. cit \unsure{lem. 2|lex. 2|cor. 2} et cor. 3}

% archive sheet number: 47
